\section{Conclusion}

We presented an evidence-bound definition of autonomous research iteration
and evaluated it in one local, preregistered campaign. Two parent-linked
equation-discovery mechanisms passed development screening and then failed
their disjoint unseen panels. The system preserved both negative results and
pivoted to a controlled research-workflow study. In that matrix, the complete
loop exceeded execute-once, recovered specified failures, reproduced every
cell, and emitted no unsupported claim under the frozen evaluator. Ablations
showed distinct roles for persistent memory, failure feedback,
preregistration, and evidence gating.

The primary result is a traceable systems artifact, not proof of general
scientific autonomy. A useful next test is an independently authored,
task-level bootstrap study in which multiple research agents must generate
and execute genuinely novel repairs under the same immutable evidence
contract. Until then, the contribution is the auditable loop, its negative
scientific history, and a reproducible controlled evaluation of how that loop
handles failure.
